\section{非被約反正準因子端点における全二次障害}

以下では $\mathbf C$ 上で考え、$\partial_{ij}=\partial_i\wedge\partial_j$ と略記する。
対象は既知の端点
\[
 Q=x_0^2\partial_{14}-\frac29x_0^2\partial_{23}
   +\frac43x_0x_1\partial_{24}
   +\left(\frac23x_0x_2+x_1^2\right)\partial_{34},
 \qquad F_Q=-\frac{2}{45}x_0^5
\]
である。$Q$ の一般階数は $4$ であり、実際の反正準因子は $5H$ である。
ここでの ``非被約'' はこの因子を指す。Poisson スキームの局所環の非被約性は、以下からは結論しない。

\subsection{全接空間、軌道、および重み}

既存の divergence-free 二次 bivector モデルにおいて、二次係数は $150$ 個ある。
Jacobi 多項式を
\[
 J(P)_{ijk}=\sum_{l=0}^4
 \bigl(P_{il}\partial_lP_{jk}+P_{jl}\partial_lP_{ki}+P_{kl}\partial_lP_{ij}\bigr)
\]
と定め、$L_Q(S)=\left.\frac{d}{d\epsilon}\right|_0J(Q+\epsilon S)$ とする。
$M=(\operatorname{div},L_Q)$ は $375\times150$ 行列であり、
$375=25+10\cdot35$ である。全線形化の厳密階数は $99$ なので、
\[
 \dim\ker M=51,\qquad
 \dim T_{[Q]}\operatorname{Pois}(\mathbf P^4)=50.
\]
$\mathfrak{sl}_5$ の軌道微分の階数は $17$ で、スカラー方向 $\mathbf C Q$ を含む。
したがって射影軌道次元は $16$、軌道を除いた接空間の次元は $34$ である。
ここで ``軌道を除いた'' とはベクトル空間の商を意味し、局所的な粗モジュライ空間の存在を仮定しない。

$g_u=\operatorname{diag}(1,u,u^2,u^3,u^4)$ とすると $(g_u)_*Q=u^5Q$ である。
重み $w_i=i$ のもと、$x^\alpha\partial_{ij}$ の生の重みは
$w_i+w_j-\sum_k\alpha_kw_k$ である。保存した $34$ 個の補空間の基底
$R_0,\ldots,R_{33}$ はそれぞれ重み斉次であり、重複度は次の通りである。
\[
\begin{array}{c|rrrrrrr}
\text{生の重み}&-1&0&1&2&3&4&5\\\hline
\text{重複度}&3&7&8&7&6&2&1\\
\text{射影正規重み}&-6&-5&-4&-3&-2&-1&0
\end{array}
\]
射影正規重みは $Q$ の重み $5$ を引いたものである。
$150$ 座標、全核、全軌道微分、補空間の全成分は
\texttt{endpoint\_full\_tangent.json} に収録した。

\subsection{全二次障害の53次元性}

$R(z)=\sum_{i=0}^{33}z_iR_i$ と置く。二次までの変形
$Q+\epsilon R+\epsilon^2S$ が Poisson となるための必要十分条件は
\[
 M(S)+(0,J(R))=0
\]
である。したがって二次障害は
\[
 \kappa_2(z)=[(0,J(R(z)))]\in\operatorname{coker}M
\]
で与えられる。これは一次変形の二次への延長障害であり、完全な形式変形環の定義方程式と同一視しない。

$M$ の保存済みの非零 $99\times99$ 小行列を $A=M_{I,K}$ とする。
$J(R(z))$ の係数行列を、最初の divergence の $25$ 行に零を補って
$B\in\operatorname{Mat}_{375\times595}(\mathbf Q)$ と書く。
列は $z_iz_j$ ($i\leq j$) に対応し、$595=34\cdot35/2$ である。
そのとき
\[
 C=B-M_{:,K}A^{-1}B_{I,:}
\]
が全二次障害の係数行列になる。
実際、$M_{:,K}$ は $\operatorname{im}M$ の基底であり、$I$ 行への射影はこの像上で同型である。
したがって $C(z_iz_j)=0$ は、$(0,J(R(z)))\in\operatorname{im}M$ と同値である。

\paragraph{厳密計算結果.}
$B$ の非零係数は $9041$ 個あり、$C$ の階数は厳密に $53$ である。
従って全 $34$ 変数上の二次障害は、線形独立な $53$ 本の二次式で完全に記述される。
独立な二次式の重み別本数は
\[
\begin{array}{c|rrrrrr}
\text{生の重みの和}&-2&-1&0&1&2&3\\\hline
\text{本数}&6&13&13&11&7&3
\end{array}
\]
となる。ここで ``独立'' は二次式の $\mathbf Q$ ベクトル空間における独立性を意味し、
$53$ 個の幾何学的に独立な制約や余次元 $53$ を意味しない。

独立性と完全性には次の二つの有理数証明書を用いた。
RREF から得た $53\times595$ 行列を $E$、その pivot 列を $K_2$ とすると、
\[
 E_{:,K_2}=I_{53},\qquad C=C_{:,K_2}E
\]
が係数ごとに成立する。これが階数の上界を与える。
さらに $C$ の保存した $53\times53$ 小行列の行列式は非零である。
独立検証器はこの行列式を整数 Bareiss 消去で再計算する。
全式は \texttt{quadratic\_obstructions.json} の
\texttt{independent\_quadratic\_equations} に、
元の射影済み行から直接選んだ独立な $53$ 式は \texttt{short\_basis} に保存した。

\subsection{最低重みの三方向と具体的な非延長例}

生の重み $-1$ の基底は $R_{19},R_{30},R_{33}$ である。
重み $-2$ の障害式 $6$ 本は、正確に
\[
 z_{19}^2,\quad z_{19}z_{30},\quad z_{19}z_{33},\quad
 z_{30}^2,\quad z_{30}z_{33},\quad z_{33}^2
\]
である。したがって $\mathbf C$ 値の一次変形が二次に延長するなら、必ず
$z_{19}=z_{30}=z_{33}=0$ である。
この三方向は Zariski 接空間には存在するが、非零のそのような成分を持つ一次変形は二次へ延長しない。
これは二次障害錐の被約な台についての主張であり、元の局所環に冪零元があることの証明ではない。

この現象には小さい証明書もある。次の bivector を取る。
\[
 R=24x_0x_4\partial_{12}+18x_1x_4\partial_{13}
      -16x_2x_4\partial_{23}+x_4^2\partial_{34}=R_{33}.
\]
直接計算で $\operatorname{div}R=0$、$L_Q(R)=0$ であり、
\[
 J(R)=72x_0x_4^2\partial_{123}.
\]
三次 trivector に対する線形汎関数
\[
 \ell(T)=[x_0x_4^2]T_{123}-\frac13[x_0x_4^2]T_{024}
\]
を取ると、全ての二次 bivector $S$ について $\ell(L_Q(S))=0$ である。
これは divergence-free 条件を課す前の全 $150$ 次元について成立する。
重みから寄与し得る $S$ は $x_4^2\partial_{01}$ のみであり、このとき
$[x_0x_4^2]L_Q(S)_{123}=4/9$、
$[x_0x_4^2]L_Q(S)_{024}=4/3$ なので主張が従う。
一方 $\ell(J(R))=72\ne0$ である。
従って $Q+\epsilon R$ は $\mathbf C[\epsilon]/(\epsilon^2)$ 上の変形だが、
$\mathbf C[\epsilon]/(\epsilon^3)$ へは延長できない。

\subsection{既知の二枝との整合性}

別途構成した $P_A,P_B$ の $19$ パラメータ族を $Q$ で微分し、軌道方向を除く。
その像を $V_A,V_B\subset\mathbf C^{34}$ とすると、厳密に
\[
 \dim V_A=\dim V_B=11,\quad
 \dim(V_A\cap V_B)=10,\quad
 \dim(V_A+V_B)=12
\]
である。さらに $53$ 本の二次障害は $V_A+V_B$ 全体で恒等的に零になる。
従ってこの二つの既知の枝は、その微分像を張る線形空間上では二次障害だけで分離されない。
これは $12$ 次元の実際の Poisson 線形族の存在を主張するものではない。
また、$V_A,V_B$ が各既約成分の全 Zariski 接空間であるとも主張しない。

\subsection{\texorpdfstring{$5H$}{5H} 族の全パラメータにわたる接次元}

$Q$ の $\partial_{14}$ の係数だけを変えた族
\[
 Q_a=ax_0^2\partial_{14}-\frac29x_0^2\partial_{23}
       +\frac43x_0x_1\partial_{24}
       +\left(\frac23x_0x_2+x_1^2\right)\partial_{34}
\]
は全て Poisson である。$a\ne0$ のとき一般階数 $4$、
$F_{Q_a}=-2ax_0^5/45$ である。
全線形化の階数と接次元は次のように完全に決定できる。
\[
\begin{array}{c|c|c|c|c}
 a&\operatorname{rank}M_a&\dim T_{[Q_a]}\operatorname{Pois}
 &\dim \operatorname{PGL}_5[Q_a]&\text{軌道を除いた接次元}\\\hline
 a\notin\{0,2/9\}&99&50&16&34\\
 a=2/9&98&51&16&35\\
 a=0&97&52&16&36
\end{array}
\]
$a=0$ は一般階数 $2$ の点であり、$5H$ の rank-$4$ 開部分から外れる。
元の端点 $Q=Q_1$ はこの族の接次元跳躍点ではない。

証明は、$M_a$ をトーラス重み別に分解し、各ブロックで
$\mathbf Q(a)$ 上の非零小行列と階数上界の行列恒等式を検証することによる。
選択した小行列の零点は
\[
 a\in\{0,1/9,2/9,1/3,8/9,-4/9,22/81\}
\]
に限られ、それぞれで別の有理数小行列と同じ上界恒等式を検証した。
実際の階数降下は $0,2/9$ のみである。
これは有限個のサンプルからの推測ではなく、全ての複素数 $a$ に対する証明である。

軌道行列にも同様の証明を適用できる。一般点の $17\times17$ 小行列の行列式は
\[
 -\frac{409600}{14348907}a^5(9a+4)
\]
であり、その零点 $a=0,-4/9$ においても代わりの小行列は非零である。
従って軌道微分の階数は全ての $a$ で $17$。
$\operatorname{diag}(-2,-1,0,1,2)$ の無限小作用が $5Q_a$ を与えるので、
射影軌道次元は常に $16$ である。

定理 \ref{thm:kappa_classification} により、五係数族では
$\kappa=AE/(CD)$ が全 $\PGL_5$ 共役の完全不変量である。
本族では $\kappa=9a/8$ なので、rank4 の唯一の接次元跳躍点は
$\kappa=1/4$ に対応する。

\subsection{検証と未決事項}

\texttt{verify\_obstructions.py} は、生成器が用いた Fraction による基底対ごとの括弧計算とは別に、
一つの記号的 $R(z)$ の Jacobi 多項式を微分・展開し、$375\times595$ 行列を再構成する。
線形化、全核、全軌道方向は標準ライブラリのみの多項式計算で再構成し、
階数下界の行列式は整数 Bareiss 消去で検証する。
パラメータ族については $\mathbf Q(a)$ 上の行列恒等式と全例外候補点を検証する。

現時点で確定したのは全二次障害であり、全高次 Kuranishi 方程式ではない。
二次障害の消滅から形式的・解析的な延長可能性を結論することはできない。
同様に、局所枝の完全列挙、局所環の被約性、二次障害錐の既約分解は未決である。
本計算は既知の全体成分の下限を増やす根拠としては用いない。
